% Осень26!B25; нумерация внутри столбца без пустых строк.
% Источник: лекция 4, §9 (строка 294)
\definitionticket{22}{Множества Лебега измеримой функции}
\label{definition:22}

\subsection*{Множества Лебега функции}

\textbf{Соглашения.} Имеется пространство с мерой \((X,\mathcal{C},\mu)\): пространство \(X\), σ-алгебра \(\mathcal{C}\subset 2^X\), мера \(\mu\) на \(\mathcal{C}\). «Мера пока не играет» --- для определений этого параграфа важна только σ-алгебра; мера понадобится позже.

Рассматриваем функцию \[f:E\to\overline{\mathbb{R}},\qquad E\in\mathcal{C},\] где \(\overline{\mathbb{R}}=\mathbb{R}\cup\{\pm\infty\}\): \textbf{очень часто будем разрешать функциям принимать бесконечные значения} («цинично говорим: значение равно \(+\infty\) или \(-\infty\)»). Требование \(E\in\mathcal{C}\) (измеримость области определения) заложено в само понятие «измерима на \(E\)».

\begin{definition}

\emph{множества Лебега.} Для \(a\in\mathbb{R}\): \[E(f<a):=\{x\in E:\ f(x)<a\},\] и аналогично с другими знаками: \[E(f\le a),\qquad E(f>a),\qquad E(f\ge a).\] Все четыре типа множеств называются \textbf{множествами Лебега} (лебеговыми множествами) функции \(f\). Параметр \(a\) подразумевается («я тычу в какое-то лебегово множество --- такое-то»).

Если \(f\) задана на всём \(X\), пишут \(X(f<a)\) и т.п.; в этом случае букву области часто опускают.

\end{definition}

\subsubsection*{Замечание 1: множества Лебега выражаются друг через друга}

Дополнение берётся \textbf{в множестве \(E\)} (естественный контекст для функций, заданных на \(E\)): \[E(f>a)=E\setminus E(f\le a),\qquad E(f\ge a)=E\setminus E(f<a).\] Очевидно: точки, где \(f>a\), --- это в точности те точки \(E\), где не выполнено \(f\le a\).

\subsubsection*{Замечание 2: нестрогие неравенства через строгие (счётные операции)}

\[E(f\le a)=\bigcap_{n=1}^{\infty}E\Bigl(f<a+\frac1n\Bigr).\] Действительно, \(f(x)\le a\iff\forall n\ f(x)<a+1/n\). «Бесхитростный» вариант \(\bigcap_{b>0}E(f<a+b)\) --- \textbf{несчётное} пересечение, «в наших науках не очень интересно»; дискретизация \(b=1/n\) даёт \textbf{счётное} пересечение --- и «в плане наблюдений за измеримостью не выводит за пределы σ-алгебры». В этом ценность замечания. Аналогично для противоположных знаков: \(E(f\ge a)=\bigcap_nE(f>a-1/n)\).
